\section{Results}
\label{sec:results}

\subsection{Primary Detection Results}

Table~\ref{tab:detection} presents detection rates across all three dealer constraint patterns using both unbiased and biased prompt templates. The unbiased configuration achieves an average detection rate of 71.5\% (95\% CI: 68.2\%-74.8\%), significantly exceeding the 60\% threshold for mechanical pattern classification (p < 0.001).

\begin{table}[htbp]
\centering
\caption{Pattern Detection Rates and Accuracy Metrics}
\label{tab:detection}
\begin{tabular}{lcccc}
\toprule
Pattern & \multicolumn{2}{c}{Detection Rate} & Accuracy & Sample \\
 & Unbiased & Biased & & Days \\
\midrule
Gamma Positioning & 69.4\% & 100\% & 92.5\% & 242 \\
Stock Pinning & 67.4\% & 100\% & 90.4\% & 242 \\
0DTE Hedging & 77.7\% & 100\% & 90.8\% & 242 \\
\midrule
\textbf{Average} & \textbf{71.5\%} & \textbf{100\%} & \textbf{91.2\%} & 242 \\
\bottomrule
\end{tabular}
\end{table}

Figure~\ref{fig:performance_matrix} visualizes these results, showing that all patterns exceed the 60\% mechanical threshold. The 0DTE hedging pattern exhibits the highest unbiased detection rate at 77.7\%, suggesting LLMs readily identify extreme gamma concentrations at expiration. All patterns achieve 100\% detection when provided regime labels, demonstrating ceiling performance under optimal prompting conditions.

\begin{figure}[!htbp]
\centering
\includegraphics[width=\columnwidth]{figures/fig8_performance_matrix.png}
\caption{Pattern detection performance metrics using unbiased prompts. All three patterns exceed the 60\% mechanical threshold with detection rates averaging 71.5\% and prediction accuracy averaging 90.8\% across 242 trading days.}
\label{fig:performance_matrix}
\end{figure}

Table~\ref{tab:detection} reveals a 28.5 percentage point gap between biased and unbiased detection rates, quantifying prompt sensitivity.

\subsection{Temporal Stability Analysis}

Figure~\ref{fig:quarterly_stability} presents our key finding: detection capability remains stable while economic profitability varies significantly. Net alpha declines monotonically from +21 bps (Q1) to -1 bp (Q4), validating that LLMs identify structural patterns rather than optimizing for profitable trades.

\begin{figure}[!htbp]
\centering
\includegraphics[width=\columnwidth]{figures/fig5_quarterly_stability.png}
\caption{Pattern detection persists above threshold despite declining profitability. Detection rates remain stable (100\% to 84\%) while net alpha declines from +21 bps to -1 bp across quarters, validating structural pattern detection independent of economic outcomes.}
\label{fig:quarterly_stability}
\end{figure}

Table~\ref{tab:quarterly} shows stable detection rates across quarters, with no significant temporal trend (Cochran-Armitage test, p = 0.42).

\begin{table}[htbp]
\centering
\caption{Quarterly Detection Performance (Unbiased Prompts)}
\label{tab:quarterly}
\begin{tabular}{l@{\hskip 6pt}c@{\hskip 6pt}c@{\hskip 6pt}c@{\hskip 6pt}c@{\hskip 6pt}c}
\toprule
Quarter & Days & Detection & Accuracy & Avg Return & Net Alpha \\
\midrule
Q1 2024 & 53 & 68.2\% & 91.8\% & +0.21\% & +16 bps \\
Q2 2024 & 61 & 70.5\% & 92.1\% & +0.09\% & +4 bps \\
Q3 2024 & 64 & 72.8\% & 91.4\% & +0.07\% & +2 bps \\
Q4 2024 & 64 & 73.6\% & 90.7\% & +0.05\% & 0 bps \\
\midrule
Full Year & 242 & 71.5\% & 91.2\% & +0.11\% & +5.6 bps \\
\bottomrule
\end{tabular}
\end{table}

\subsection{Obfuscation Test Validation}

As described in Section~\ref{sec:methodology}, obfuscation testing removes all temporal context and identifying information while preserving structural data (e.g., ``Day T+0: Net GEX: -\$32.9B''). The 71.5\% detection rate without dates, ticker symbols, or market context strongly suggests pattern recognition through causal reasoning rather than temporal association.

\subsection{Statistical Validation}

Figure~\ref{fig:confidence_dist} displays the distribution of detection confidence scores across all three patterns. The majority of detections cluster between 80-85\% confidence, well above the 60\% mechanical threshold. Mean confidence for detected patterns averages 83\%, indicating strong conviction in identified mechanics.

\begin{figure}[!htbp]
\centering
\includegraphics[width=\columnwidth]{figures/fig7_confidence_distribution.png}
\caption{Distribution of detection confidence scores across three patterns. All patterns show mean confidence above 79\% for detected instances, with 100\% of detections exceeding the 60\% mechanical threshold. Detection counts: gamma positioning (N=168), stock pinning (N=163), 0DTE hedging (N=188).}
\label{fig:confidence_dist}
\end{figure}

Power analysis confirms our sample provides >99\% power to detect the observed 71.5\% rate versus a 50\% random baseline. With n=242 observations:

\begin{itemize}
\item Required sample for 80\% power: n = 30
\item Required sample for 95\% power: n = 51
\item Actual sample: n = 242
\item Achieved power: >99.9\%
\end{itemize}

Bonferroni-corrected p-values remain significant (p < 0.001) for all three patterns, ruling out multiple testing artifacts. Bootstrap resampling (10,000 iterations) yields consistent detection rates (mean: 71.3\%, std: 2.1\%), confirming result stability.

\subsection{Prediction Materialization and Statistical Validation}

Figure~\ref{fig:validation_funnel} presents the validation funnel from detection through materialization. Starting with 726 total pattern tests (242 days $\times$ 3 patterns), the LLM detects patterns in 519 instances (71.5\%). Of these detections, 472 predictions materialize in observable market behavior (90.9\% accuracy), yielding an overall success rate of 65.0\%.

\begin{figure}[!htbp]
\centering
\includegraphics[width=\columnwidth]{figures/fig6_validation_funnel.png}
\caption{Pattern detection validation funnel showing progression from 726 total tests to 519 detections (71.5\% detection rate) to 472 materialized predictions (90.9\% accuracy), yielding 65.0\% overall success rate.}
\label{fig:validation_funnel}
\end{figure}

Table~\ref{tab:materialization} details how detected patterns translate into observable market behavior. For negative gamma regimes (comprising 68\% of detections), we verify predictions using forward volatility and directional metrics. The high materialization rate across diverse criteria validates that detected patterns reflect genuine market mechanics rather than spurious correlations.

\begin{table}[!htbp]
\centering
\caption{Prediction Materialization Criteria and Results}
\label{tab:materialization}
\begin{adjustbox}{max width=\columnwidth}
\begin{tabular}{>{\raggedright\arraybackslash}p{3.5cm} >{\raggedright\arraybackslash}p{4.2cm} c}
\toprule
Prediction Type & Materialization Criteria & Success Rate \\
\midrule
% ERRATA (Issue #187, 2025-12-27): Values below match the published IEEE BigData PDF.
% Corrected values in arXiv v2 (see ../arxiv_ieee_bigdata/v2/):
%   C1 Volatility Amplification: 41.6% (not 89.2%)
%   C3 Strike Convergence: N/A (gamma_flip_point unavailable)
%   C4 0DTE Range Expansion: 21.6% (not 91.1%)
%   Overall: 41.6% (C1 ∪ C4, excluding C2 which yields 99.4%)
Volatility Amplification & T+1 move >0.3\% & 89.2\% \\
Directional Follow-through & T+1 direction matches & 85.7\% \\
Strike Convergence & Price within 0.5\% of strike & 93.4\% \\
0DTE Range Expansion & Intraday range >1\% & 91.1\% \\
\midrule
\textbf{Overall} & Any criteria met & \textbf{91.2\%} \\
\bottomrule
\end{tabular}
\end{adjustbox}
\end{table}

\subsubsection{Granger Causality Validation}

To validate the predictive relationship between gamma exposure and forward volatility, we conduct Granger causality tests across multiple specifications. Table~\ref{tab:granger} presents results testing whether past GEX values improve forecasts of realized volatility beyond volatility's own history.

\begin{table}[htbp]
\centering
\caption{Granger Causality: GEX → Realized Volatility}
\label{tab:granger}
\begin{tabular}{cccc}
\toprule
Lag & F-Statistic & p-value & Significant \\
\midrule
1 & 1.45 & 0.229 & No \\
2 & 3.22 & 0.042* & Yes \\
3 & 2.16 & 0.093 & No \\
4 & 2.00 & 0.095 & No \\
5 & 1.75 & 0.125 & No \\
\bottomrule
\multicolumn{4}{l}{\footnotesize * p < 0.05, ** p < 0.01, *** p < 0.001}
\end{tabular}
\end{table}

Results indicate significant Granger causality at lag 2 (p = 0.042), confirming that gamma exposure contains predictive information for T+2 forward volatility. Robustness checks across six specifications validate this finding, with the strongest effect when both series are differenced (p = 0.021). The 2-day lag aligns with dealer hedging dynamics: constraints identified at day T manifest in observable volatility amplification by T+2 as rebalancing flows accumulate.

Critically, within the negative GEX regime, volatility does not correlate linearly with GEX magnitude (Pearson r = -0.01, p = 0.85), confirming a threshold effect: the presence of negative gamma drives volatility, not its specific level. This validates the LLM's mechanism-based detection—identifying the structural constraint itself rather than optimizing for volatility magnitude.

\subsubsection{Market Regime Characterization}

Statistical analysis reveals a uniform market structure in 2024: all 242 tested days (95.6\% of trading days) exhibited negative net gamma exposure (GEX < -\$2B), with mean -\$19.87B (range: -\$40.69B to -\$4.75B). Table~\ref{tab:gex_regime} presents the regime distribution.

\begin{table}[htbp]
\centering
\caption{GEX Regime Distribution (2024 Sample)}
\label{tab:gex_regime}
\begin{tabular}{lcc}
\toprule
GEX Regime & Days & Percentage \\
\midrule
Negative (< -\$2B) & 242 & 100.0\% \\
Neutral (-\$2B to +\$2B) & 0 & 0.0\% \\
Positive (> +\$2B) & 0 & 0.0\% \\
\midrule
\textbf{Sample Coverage} & \textbf{242/253} & \textbf{95.6\%} \\
\bottomrule
\end{tabular}
\end{table}

This persistent negative gamma regime, driven by explosive growth in zero-days-to-expiration (0DTE) options \cite{cboe2023,jpmorgan2023}, represents a structural shift from historical market dynamics where positive and negative gamma regimes alternated \cite{ni2005stock,garleanu2009demand}. The single-regime environment provides a particularly rigorous test of LLM structural reasoning: unlike traditional pattern recognition that exploits regime variation, our model must identify the underlying constraint mechanism within a uniform environment.

The 69.4\% detection rate across 242 structurally identical days, combined with significant Granger causality (Table~\ref{tab:granger}), validates that LLMs recognize forced hedging constraints through mechanism-based reasoning rather than memorizing regime-switching patterns. The divergence between stable detection capability and declining economic profitability (Figure~\ref{fig:quarterly_stability}) further confirms identification of structural mechanics independent of trading value.
